\DocumentMetadata{
  pdfversion=2.0,
  pdfstandard=ua-2,
  lang=en-US,
  tagging=on,
  tagging-setup={math/setup=mathml-SE,tabsorder=structure}
}
\documentclass[11pt,letterpaper]{article}
\usepackage[margin=0.68in,includefoot]{geometry}
\usepackage{newcomputermodern}
\usepackage{amsmath,amscd,mathtools}
\usepackage{tikz,adjustbox}
\providecommand{\square}{\mdlgwhtsquare}
\usepackage[hidelinks]{hyperref}
\hypersetup{
  pdftitle={Algebra First-Year Exam, August 2018, Part 2},
  pdfauthor={University of Florida Department of Mathematics},
  pdfsubject={Graduate mathematics examination},
  pdfdisplaydoctitle=true,
  bookmarksopen=true
}
\setcounter{secnumdepth}{0}
\setlength{\parindent}{0pt}
\setlength{\parskip}{0.28em}
\setlength{\emergencystretch}{3em}
\setlength{\leftmargini}{2.15em}
\setlength{\leftmarginii}{2.65em}
\setlength{\leftmarginiii}{3em}
\renewcommand{\labelenumi}{\textbf{\theenumi.}}
\pagestyle{empty}
\raggedbottom
\allowdisplaybreaks
\newenvironment{examproblems}[1]{%
  \begin{enumerate}%
  \setcounter{enumi}{#1}%
  \setlength{\topsep}{0.35em}%
  \setlength{\partopsep}{0pt}%
  \setlength{\itemsep}{0.48em}%
  \setlength{\parsep}{0.18em}%
}{\end{enumerate}}
\newenvironment{examsubparts}{%
  \begin{enumerate}%
  \setlength{\topsep}{0.18em}%
  \setlength{\partopsep}{0pt}%
  \setlength{\itemsep}{0.16em}%
  \setlength{\parsep}{0.1em}%
}{\end{enumerate}}
\newenvironment{examinstructions}{%
  \par\begingroup\itshape
}{%
  \par\endgroup\vspace{0.18em}
}
\makeatletter
\renewcommand\section{\@startsection{section}{1}{0pt}%
  {-1.25ex plus -.25ex minus -.1ex}%
  {0.55ex plus .1ex}%
  {\normalfont\large\bfseries}}
\renewcommand{\@maketitle}{%
  \newpage\null\vskip -1em%
  {\footnotesize\bfseries
    \href{https://www.ufl.edu/}{University of Florida}, \href{https://math.ufl.edu/}{Department of Mathematics}\par}%
  \vskip -0.5em%
  \tagstructbegin{tag=Title}%
    {\LARGE\bfseries\href{https://gma.math.ufl.edu/exams/algebra-first-year/}{Algebra First-Year Exam}, August 2018, Part 2\par}%
  \tagstructend%
  \vskip 0.25em%
  \tagmcbegin{artifact}%
    \hrule height 1pt%
  \tagmcend%
  \vskip 1em%
}
\makeatother
\title{Algebra First-Year Exam, August 2018, Part 2}
\author{}
\date{}
\begin{document}
\maketitle
\thispagestyle{empty}
\begin{examinstructions}
Answer four problems. Write your solutions clearly in complete English sentences. You may quote results (within reason) as long as you state them clearly.
\end{examinstructions}
\begin{examproblems}{0}
\item
Let~\(R\) be a commutative ring with~\(1\) and let \(I, J, P\) be ideals in~\(R\) such that \(I \cap J \subset P\).
\begin{examsubparts}
\item[\textnormal{(a)}]
Prove that if~\(P\) is a prime ideal then either \(I \subset P\) or \(J \subset P\).
\item[\textnormal{(b)}]
Give an example which shows that if~\(P\) is not prime then the conclusion above may not hold.
\end{examsubparts}
\item
Prove that the ring \(\mathbb{Z}[\sqrt{-2}] = \{a + b\sqrt{-2} : a,b \in \mathbb{Z}\}\) is a Euclidean domain with respect to the norm defined by \(N(a+b\sqrt{-2})=a^2+2b^2\) for \(a,b \in \mathbb{Z}\).
\item
Prove that the following polynomials are irreducible over~\(\mathbb{Q}\). Cite any theorems that you use.
\begin{examsubparts}
\item[\textnormal{(a)}]
\(X^3-5X^2+X+2\)
\item[\textnormal{(b)}]
\(X^4+33X^2-6\)
\item[\textnormal{(c)}]
\(X^4+2X^3-14X^2+31X-7\)
\end{examsubparts}
\item
Let~\(R\) be a ring with~\(1\).
\begin{examsubparts}
\item[\textnormal{(a)}]
State the First Isomorphism Theorem for~\(R\)-modules.
\item[\textnormal{(b)}]
Let \(A_1,A_2,\ldots,A_n\) be~\(R\)-modules and for \(1 \leq i \leq n\) let \(B_i\) be a submodule of \(A_i\). Prove that 
\[
(A_1 \times \cdots \times A_n)/(B_1 \times \cdots \times B_n) \cong (A_1/B_1) \times \cdots \times (A_n/B_n).
\]
\end{examsubparts}
\item
State the classification theorems for finitely generated modules over a PID. You should give both the invariant factor and the elementary divisor decompositions. Be sure to include the uniqueness statements.
\item
Let \(L/K\) be an algebraic field extension and let~\(R\) be a subring of~\(L\) such that \(K \subset R\). Prove that~\(R\) is a field.
\end{examproblems}
\end{document}
